\scanpage{8}
% archive sheet number: 7

On sait que ($X_0\longrightarrow Y_0$ \unsure{étant} couvrt qu. cpct qu. sép.)

\[
\check H^1(X_0/Y_0,\underline{\unsure{I|J|\mathcal J}}_{Y_0})=0
\]

donc on peut résoudre l'équation $d(\lambda)=y$, ce qui prouve la surjectivité annoncée.

\deleted{Soit} Soit $Y=\operatorname{Spec} A$, alors \deleted{$A$} on obtient un diagramme commutatif

\diagram{\begin{tikzcd}[ampersand replacement=\&]
Y\&X\arrow[l]\&R\arrow[l,shift left=.8ex]\arrow[l,shift right=.8ex]\\
Y_0\arrow[u,<->]\&X_0\arrow[l]\arrow[u]\&R_0\arrow[l,shift left=.8ex]\arrow[l,shift right=.8ex]\arrow[u]
\end{tikzcd}}

D'ailleurs, \deleted{\unsure{posons|pour} $\underline I_{Y_0}$} par l'exactitude écrite précédemment, $Y_0$ est un sous-schéma fermé de $Y$ défini par un idéal nilpotent isomorphe à $\underline I_{Y_0}$.

On voit de suite que

\[
\underline I_{Y_0}X=\underline I_{X_0}\simeq\underline I_{Y_0}\unsure{\otimes_{Y_0}|\mathcal O_{Y_0}|\otimes_{\mathcal O_{Y_0}}}\mathcal O_{X_0}
\]

comme Idéal sur $X$, i.e. $X_0\longrightarrow Y_0$ déduit de $X\longrightarrow Y$ par changement de base $Y_0\longrightarrow Y$. Comme de plus $X_0/Y_0$ est plat, $X/Y$ est \uline{plat}. Comme $R\longrightarrow X$ est plat, $R$ est plat sur $Y$, et $R\longrightarrow X\times_Y X=P$ est un morphisme de $Y$-préschémas \uline{plats}, tels que $R_0\longrightarrow P_0$ est un isom, donc $R\longrightarrow P$ est un isomorphisme, O.K.
